%! root=./main.tex
\section{Introduction}

{ % new; CCS and CFI-oriented.
  % NOTE (@kjee): 1 page + 1/3 ~ 2/3 pages with about 3 contribution bullets.
  %%%%%%%%%%%%%%%%%%%%%%%%%%%%%%%%%%%%%%%%%%%%%%%%%%%%%%%%%%%%%%%%%%%%%%%%%%%%%%%%
  % CFI is the king of modern security defense.
  %%%%%%%%%%%%%%%%%%%%%%%%%%%%%%%%%%%%%%%%%%%%%%%%%%%%%%%%%%%%%%%%%%%%%%%%%%%%%%%%
  \kjee{we may want to have a sentencer that explains \cfi.}
  CFI has become a cornerstone defense against control-hijacking attacks, enabling widespread deployment with strong guarantees and near-native overhead. Yet the community focuses on design-level security guarantees, runtime overheads, and policy generation aspects, while largely overlooking the correctness of CFI implementations in production, leaving a critical gap between theory and practice. This gap underscores a broader challenge in security engineering: translating sound designs into correct, deployable systems is difficult and costly. Even rigorously designed systems often contain subtle flaws that emerge only after deployment. In security-critical settings, these flaws do more than degrade functionality—they undermine the trusted computing base, create false confidence, and expose exploitable weaknesses.
  %
  To mitigate these risks, the community relies primarily on dynamic testing techniques such as fuzzing and symbolic execution. While effective at uncovering concrete bugs, these approaches remain inherently limited. This limitation stems not from engineering shortcomings, but from the fundamental nature of dynamic analysis.

  %%%%%%%%%%%%%%%%%%%%%%%%%%%%%%%%%%%%%%%%%%%%%%%%%%%%%%%%%%%%%%%%%%%%%%%%%%%%%%%%%%%%%%%%%%%%%%
  % Now FM story begins; FM to the rescue.
  %%%%%%%%%%%%%%%%%%%%%%%%%%%%%%%%%%%%%%%%%%%%%%%%%%%%%%%%%%%%%%%%%%%%%%%%%%%%%%%%%%%%%%%%%%%%%%

  % TODO (@kjee): In FM (or Coq) ecosystem, I think we have two different types of sources, say (1) *.v and (2) *_proof.v. Can we name them and call them in a specific way?
  % TODO (@kjee): re-work on sentence ordering.

  Formal methods take a fundamentally different approach. Rather than observing executions, formal methods establish stringent correctness guarantees by proving that an implementation satisfies formally specified properties, thereby covering all program behaviors within the scope of the specification. Proof assistants then mechanically check the resulting proof artifacts, reducing trust to the correctness of the specification and the soundness of the proof checker.

  Despite these unparalleled advantages, formal methods have seen limited adoption in production systems for three reasons:
  % NOTE (kjee): the following is in a way known; conventional perception... However, let's not think about/mention it here at least.
  % Formal verification requires substantial effort and specialized expertise to construct proof artifacts, often rivaling or exceeding the cost of the implementation itself. Moreover, these proofs are tightly coupled to the verified code: even minor changes --- such as bug fixes, optimizations, or refactoring --- can invalidate existing proofs and trigger significant re-verification effort. This tight coupling makes it costly to maintain proofs alongside evolving codebases and limits scalability to large or heterogeneous systems.
  First, the guarantees formal methods deliver are only as strong as the specifications themselves, which may be incomplete or flawed.
  Second, proofs bind tightly to specific implementations and verification contexts, making it difficult to scale verification across large, heterogeneous, and evolving codebases.
  % TODO (@kjee): my first attempt to justify/motivate Picinae layer/framework.
  Third, formally verifying software at the binary level has remained out of reach because the community has lacked a sound and accurate abstraction layer to serve as a foundation for such verification.

  %%%%%%%%%%%%%%%%%%%%%%%%%%%%%%%%%%%%%%%%%%%%%%%%%%%%%%%%%%%%%%%%%%%%%%%%%%%%%%%%%%%%%%%%%%%%%%
  % Our story begins; formally-verified, machine-checked CFI.
  %
  %  - In this research, we present the first formally-verified (or machine-check) CFI system that extends static binary writer framework, built for fully mature, realistic architecture.
  %    - in an effort to impending, imminent, unanswered questions for the safety and security of a critical security defense where compliance to the design specification and runtime correctness still remain at the hands of dynamic testing which is known to be limited and incomplete.
  %%%%%%%%%%%%%%%%%%%%%%%%%%%%%%%%%%%%%%%%%%%%%%%%%%%%%%%%%%%%%%%%%%%%%%%%%%%%%%%%%%%%%%%%%%%%%%

  % TODO (@kjee): we need a killer setence here; summarizing two-folds (CFI implementation accuracy concern, FM scalability on arbitary binaries inputs) as a connecting/strong/motivational statement.
  To respond, our research takes a step toward closing that gap. We present the first machine-checked CFI system built on a static binary rewriting framework for a mature, real-world architecture, providing formally verified compliance with the design specification and runtime correctness beyond the incomplete assurances of dynamic testing.

  % NOTE (@kjee): specificiation to formally verifiying CFI can be straightforward.
{
    % Binary CFI is a well-established and well-specified defense.
    \pname implements source-agnostic (binary-only) CFI by rewriting binaries to confine all control-flow transfers to policy-specified edges. It identifies and instruments control-transfer instructions (\eg jumps, calls, and returns), checking each target against the policy at runtime and triggering a safe abort on violation.
    By definition,\wc{formally} the CFI enforcement for a given binary is \emph{sound} if all possible execution traces of every rewritten binary obey the control-flow policy.
}

  % NOTE (@kjee): FM of binary writer itself is a difficult, unexplored, open challenge.
\comment{
- Binary-level CFI, built on static binary rewriting techniques, is a well-studied and mature area.
  - However, formal verification of CFI implementation correctness remains understudied, despite the long-standing availability of mature proof assistants (Coq/Rocq, Isabelle/HOL).
  - The fundamental challenge is the absence of ISA abstractions that are both faithful to real-world architectures and amenable to formal reasoning.
  - As a result, prior work has been largely confined to reasoning about concrete binaries known ahead of time, and cannot generalize to arbitrary binaries supplied by users.
    - For instance, RockSalt successfully verified the safety guarantees of an SFI solution for x86, but its analysis is confined to a concrete subset of the x86 ISA and the specific binaries conforming to NaCl's sandbox policy, not generalizeable to arbitrary binaries or other ISAs.

- We address this challenge by building on the Picinae framework~\cite{picinae}, which provides a principled foundation for machine-checked proofs about real binary programs on real ISAs. Critically, Picinae supports proofs in which the target binary is treated as a symbolic proof variable rather than a fixed program, enabling verification that generalizes across arbitrary binaries.
}

Binary-level CFI, built on static binary rewriting techniques, is a well-studied and mature area. Yet formally verifying the correctness of a binary CFI implementation, let alone the broader class of binary rewriters, remains an open challenge, despite the long-standing availability of mature proof assistants such as Coq/Rocq and Isabelle/HOL. The core difficulty is one of semantic generality: correctness must hold for all possible input binaries, whether compiler-generated, hand-written, or otherwise, not merely those drawn from a particular toolchain or restricted to a subset of the ISA. Prior work has been largely confined to reasoning about concrete, fixed binaries, falling short of this bar. RockSalt, for instance, successfully verified the safety guarantees of an SFI solution for x86, but its analysis is tied to a concrete subset of the x86 ISA and the NaCl sandbox policy, and does not extend to arbitrary binaries or other ISAs.

\optional{%
  Formal verification directly targets this challenge by reasoning over all possible inputs rather than sampled executions, enabling universal correctness guarantees that dynamic testing cannot provide. These guarantees hold regardless of binary provenance, compiler, or ISA quirk, making formal methods particularly well suited to establishing that a CFI implementation soundly enforces its policy across the full space of binaries it may encounter.
}

We address the key obstacle, namely the absence of ISA abstractions amenable to formal reasoning, by building on the Picinae framework~\cite{picinae}, which supports machine-checked proofs in which the target binary is treated as a symbolic proof variable rather than a fixed program. This enables \pname to achieve verification that generalizes across arbitrary binaries, a capability no prior binary CFI rewriter has demonstrated.

\comment{%
  Yet formally verifying a binary CFI, let alone the broader class of binary rewriter implementations, remains an open challenge. The core difficulty lies in semantic generality: correctness requires enforcing the desired safety property for all possible input binaries, not merely those produced by a particular compiler or restricted to a subset of the ISA.
  %
  Formal verification of static CFI instrumentation directly targets this challenge. By reasoning over all possible inputs rather than sampled executions, formal methods can establish universal correctness guarantees that dynamic testing cannot provide. This makes them particularly well suited to demonstrate that a CFI implementation soundly enforces its policy across the full space of binaries it may encounter.
}

\comment{% from the original intro text; the following makes sense.
  Since source-driven formal methods are not well suited to express or verify source-agnostic soundness properties at the binary level, and binary-driven approaches typically focus on verifying concrete binaries rather than \emph{the arbitrary} outputs of a transformation, prior efforts to verify binary CFI have largely resorted to case-by-case verification of CFI-hardened binaries.
  %
  This approach enables code consumers to protect themselves from untrusted producers, but it does not equip producers with a way to construct a CFI implementation that is provably sound by design and guaranteed to pass such checks. Instead, producers must rely on unit testing—applying the rewriter to many inputs and checking for violations, thereby forfeiting the full-coverage guarantees of formal methods and risking deployment-time failures when the rewriter produces an unsound binary.

  Our research addresses the challenge by incorporating Picinae framework~\cite{picianae} serves as a core enabler in addressing the intrinsic challenges of formally verifying a binary rewriter. Picinae provides a ground-up foundation for machine-checked proofs about real binary programs on real ISAs, while maintaining sufficiently general foundations to prove theorems in which all or part of the target binary code is a proof variable rather than a specific program.
}


Leveraging the framework level support,  \pname accurately models the intricate features of ARMv7, including conditional instruction prefixes, undefined and non-deterministic effects, memory access permissions, hardware exceptions, and indirect calls and returns with mixed endianness. This faithful modeling of a mature, real-world ISA provided by Picinae, \pname implements the first formally verified CFI framework for ARMv7 binaries.

Together, these design choices establish the formal verification of a CFI-enforcing static binary rewriter as a rare combination of tractability and impact. To the best of our knowledge, no prior work has explored this space, despite its direct implications for closing the long-standing assurance gap in deployed CFI systems.
%
%%%%%%%%%%%%%%%%%%%%%%%%%%%%%%%%%%%%%%%%%%%%%%%%%%%%%%%%%%%%%%%%%%%%%%%%%%%%%%%%
% specific story begins;
%%%%%%%%%%%%%%%%%%%%%%%%%%%%%%%%%%%%%%%%%%%%%%%%%%%%%%%%%%%%%%%%%%%%%%%%%%%%%%%%
% NOTE (@kjee): 2nd pass for the bullet points.
As a result, in this research, we bring in the following contributions:
\begin{itemize}[topsep=0pt,itemsep=-1ex,partopsep=1ex,parsep=1ex, leftmargin=*]
    \item  We present \pname, a machine-checked binary CFI framework that formalizes and verifies the correctness of CFI rewriting. We implement it as the first machine-checked CFI rewriter for ARMv7 and achieve near-native overhead (XX~\% average on SPEC 2017).
    % We introduce \pname, a machine-checked binary CFI framework that formally defines the security and safety conditions required for a correct CFI rewriter and verifies them against a real-world implementation. \pname is the first machine-checked binary CFI rewriter for the ARMv7 ISA, demonstrating the generality of our approach while incurring near-native runtime overhead (XX~\% average on SPEC 2017).

    \item Built on Picinae, \pname analyzes arbitrary binaries, overcoming prior FM limitations that require concrete inputs and enabling verification over unbounded binary spaces.
    % Built on the Picinae framework, \pname accepts arbitrary binaries as input, overcoming a long-standing limitation of prior FM approaches that required concrete byte sequences rather than symbolic representations,  extending the reach of formal verification to unbounded binary inputs.

    \item  We explore a previously uncharted intersection between FM and binary-level security, deriving principled FM-aware design guidelines and uncovering subtle but critical flaws that dynamic testing alone would have missed.
\end{itemize}

  % NOTE (@kjee): put the following as a closing/visionary statement (rather than one of bullet points) since it hasn't yet been realized with actual implementations.
  Looking ahead, by building a binary rewriter framework on top of the ISA abstraction layer provided by Picinae, this work takes a concrete step toward a broader and largely unexplored research direction: the formal verification of large families of security-sensitive binary rewriting tasks beyond CFI. We identify the key challenges and open problems that must be addressed to realize this vision, laying the groundwork for future research in this space.
}

\comment{% original.
  % we want to motivate the reader so that they agree that formal verification of an implementation of binary rewriting is the best way forward
  % we need to convince them that people have tried the workarounds and that they just don't provide the same level of assurance that we will

  Development of mission-critical software systems, such as secure OSes, VMs, hypervisors, and runtime systems intended for deployment in adversarial environments, demands exceptionally high-assurance development tools and procedures to meet strict safety, security, and correctness standards.
  Machine-checked formal verification has emerged as a hallmark approach for achieving the requisite guarantees for such systems,
  with success stories such as CompCert~\cite{compcert}, seL4~\cite{sel4}, and CakeML~\cite{cakeml} demonstrating the efficacy of such \emph{formal methods} (FM) for obtaining real-world, high-assurance systems in practice.
  These successes rely on mature theorem proving environments such as Rocq~\cite{rocq} and Isabelle/HOL, which have enabled rigorous, machine-checked reasoning about source-level semantics and compiler correctness.

  % \emph{Binary rewriters} transform binary programs to optimize, instrument, or enforce a security policy on them.
  % Rewriters are difficult to verify due to their complexity and the lack of tooling capable of formally reasoning about underspecified pieces of code.
  % The accepted solution to make a rewriter trustworthy is to split the task of creating a sound rewritten binary across two pieces of software.
  % The untrusted rewriter rewrites the target program without guaranteeing its output is a sound transformation.
  % The formally verified \emph{verifier} checks the output and accepts it if it is correct.

  While machine-certified compilation has received extensive attention in the literature (cf., \cite{lustre_compiler,certi_crypto_comp,certicoq_wasm,compcerto,rocq_compilers_alacarte,design_impl_certifying_compiler,certified_choreographies}), binaries that execute on real mission-critical systems are seldom the unaltered output of a compiler.
  Compiled binaries undergo a variety of low-level transformations prior to execution, including load-time linking, hardening, and sandboxing, which prepares the binary to safely execute in the context of the surrounding OS or VM\@.
  The \emph{binary rewriters} that implement these transformations demand an equivalent or higher level of assurance, since any flaw could undermine all assurances achieved during compilation.
  % NOTE (@kjee): yet we are not doing *complete* formal verification --- we only deal with .text section and rest of the parts such as (1) instruction decoder, elf archiver/relocator (?) are outside scope of the current work. *complete* may sound like an over-claim.
  However, complete formal verification of a binary rewriter \emph{implementation} has remained a difficult challenge for FM, primarily because their correctness properties are universal quantification over arbitrary binary code.
  In particular, a binary rewriter is correct if and only if it achieves a desired safety property for all possible input binaries, not merely those produced by a particular compiler or that limit themselves to a particular subset of the ISA\@.

  Binary \emph{control-flow integrity} (CFI) is a well-studied example of the associated challenges~\cite{safedispatch,typearmor,object_flow_integrity,fineibt,renewablejitcfi,eccut,resource_constrained_cfi,fixingfineibt}.
  Source-agnostic CFI algorithms transform input binaries so as to confine all their reachable control-flows to a policy-specified whitelist of edges.
  Implementations typically identify and instrument the target binary's control-flow transfer instructions (e.g., jumps, calls, and returns) with extra \emph{guard code} that checks each impending destination against the policy at runtime.
  Impending violations solicit a safe intervention, such as secure abort, rollback, or restart.
  A binary CFI implementation is \emph{sound} if all reachable execution traces of all binaries it generates obey the policy.

  Since source-driven (\emph{top-down}) FM approaches cannot state or verify a source-agnostic soundness property, and since binary-driven (\emph{bottom-up}) FM approaches are typically designed to verify concrete input binaries (not the arbitrary output of a binary transformation), prior efforts to verify binary CFI have resorted to verifying each CFI-hardened binary case-by-case (e.g., \cite{rocksalt}).
  This helps code-consumers protect themselves from untrusted code-producers, but it does not empower code-producers to create a verifiably sound CFI implementation that always passes the consumer's checks.
  Producers must resort to unit testing (i.e., running the rewriter on many input binaries to see whether any resulting binaries fail verification), losing the complete-coverage benefits of FM and risking denial of service if the deployed rewriter fails.

  Formally verifying a binary rewriter implementation requires overcoming two major hurdles:
  First, it requires a verification framework that supports \emph{underconstrained-code properties} --- properties in which the target code is a quantified proof variable, not a concrete program.
  Such properties are inexpressible in most existing bottom-up FM frameworks (e.g., \cite{trabin,isabil,islaris,myreen08,bjorn14}).
  Second, the approach must have a capability for reasoning about arbitrary byte sequences as (possibly invalid) binary code.
  This is necessary because rewriter implementations read and write bytes, and their soundness is therefore contingent upon their correct interpretation of those bytes as instructions.
  Idealizing a rewriter as an abstract model of its algorithm waives verification of many of the most complex parts of its implementation, which are pernicious sources of potential error.

  %%% KWH Comment: The relationship of this work to ours is too complex for the intro, imo.  We should discuss this in more depth in the related work section.
  %The only work in the literature which may support this was released in 2025 and focused on relational properties of \emph{concrete} programs with statically known code~\cite{relational_realistic}.


  % Verifying binary rewriters requires reasoning about their effects on arbitrary input codes.
  % Historically, binary analysis frameworks have been limited to analyzing concrete code~\cite{bap,isabil,islaris,myreen08,bjorn14,trabin}, preventing them from reasoning about binary transformations.
  % In 2025 a new binary analysis platform debuted with this capability but did not demonstrate its application, focusing instead on relational logic~\cite{relational_realistic}.

  % \emph{Top-down FM} yields high assurance for software representable as source code by extending proofs of source-level properties down to executable codes, such as through type-preserving certified compilation.
  % But codes that are binary-only, and proofs of binary-level properties that are independent of sources, require a different approach.
  % For example, low-level runtime ABIs often contain raw assembly or machine code that has been carefully optimized for maximal performance on specific instruction set architectures (ISAs).
  % These optimizations benefit all source codes that link to them, but potentially undermine top-down FM guarantees that only prove properties of the compiler-generated parts of the final executable, not the ABI code contributed by the runtime system during linking.
  % Likewise, the correctness of software that manipulates binary code, such as VMs, process loaders, linkers, and binary security hardening transforms, cannot be expressed purely as source-level properties because they concern software that takes \emph{arbitrary} binary code as an \emph{input}.
  % For example, a VM implementation is secure only if it can successfully load and execute arbitrary byte sequences (even those not generated by any compiler) in a way that prevents the loaded program from escaping the VM sandbox.

  % \emph{Bottom-up FM} verifies binary code by lifting~\cite{formally_verified_lifting} it to a low-level intermediate language (IL) that expresses the semantics of machine instructions as cpu state transitions.
  % When the IL has a formally defined semantics, it is amenable to interactive theorem proving, facilitating the same style of compositional, mechanized reasoning used in top-down FM, but applied directly to executable code.
  % This capability affords stating and proving properties of all possible executions.

  % But while existing bottom-up FM frameworks facilitate reasoning about concrete binary codes, the task of formally reasoning about arbitrary binary codes and the programs that read them has remained an unsolved challenge.
  % As a hallmark example, \emph{binary control-flow integrity} (CFI) algorithms~\textcolor{red}{[citations]} instrument input binaries to constrain their control-flows to a whitelist of permissible control-flow transfer edges.
  % The output of the CFI transform is a self-monitoring, self-securing binary that requires no additional sandboxing.
  % CFI security is a core technology for enforcing many critical software security policies, including process isolation, software fault isolation, stack integrity, mobile code security, and code-reuse attack detection and defense~\textcolor{red}{[citations throughout this list]}.

  % Although prior efforts have verified idealized binary CFI \emph{algorithms} for abstract ISA models of varying detail~\textcolor{red}{[citations]}, % no formally verified binary CFI \emph{implementation} for a real, mainstream ISA has yet been realized.
  % none have yet verified a binary CFI \emph{implementation} for a real mainstream ISA.
  % This has left open the dangerous possibility that binary CFI implementations might not match the algorithms on which their designs are informally based.
  % In practice, such implementations have many subtle details that are prone to error, including accurate decoding of all reachable instructions in potentially malicious or malformed binaries, correct identification of all possible indirect control-flow transfer operations, complete mediation of all such operations, and careful rebasing and relinking of all code fragments displaced by guard code introduced during the rewriting process.
  % A full proof of safety for a binary CFI implementation must therefore consider all (arbitrary) input byte sequences that constitute executable code on the ISA, and prove that all possible executions of the implementation's output byte sequence satisfy the control-flow policy.

  %Although checkers provide assurance when they accept a binary, they offer no guarantees or guidance when they reject one, leaving the user uncertain about the safety of the transformed code.

  In this paper we address this long-standing challenge through the design and implementation of \Picinae{}, a \uline{p}latform \uline{i}n Ro\uline{c}q for \uline{in}struction-level \uline{a}nalysis of \uline{e}xecutables.
  \Picinae{} is designed from the ground up to construct machine-checked proofs about real binary programs for real ISAs, while maintaining sufficiently general foundations to prove theorems in which all or part of the target binary code is a proof variable rather than a specific program.
  It supports modeling both von Neumann and Harvard architectures, including difficult ISA features like misaligned instructions on Intel, conditional instruction prefixes on ARM, undefined and non-deterministic effects, memory access permissions, self-modifying code, hardware exceptions, indirect calls and returns, mixed endianness on ARM, instructions with internal microcode loops, and multi-ISA integrated systems.
  These capabilities are leveraged to prove that a realistic binary CFI implementation for a mainstream ISA (ARM7) safely hardens arbitrary binaries (including rejecting unsafe or malformed binaries that it is designed to exclude from execution).
  Specifically, we present the following contributions:
  \begin{itemize}
      \item the \Picinae{} framework, which facilitates machine-checked proofs of correctness for binary programs and binary transformations;
      \item a dependently-typed, co-inductive foundation for proving correctness properties of algorithms that manipulate machine code; and
      \item a case study demonstrating the first machine-checked proof of soundness for a binary CFI rewriter implementation targeting ARM7 microprocessor code.
  \end{itemize}
}